\section{Revenue}
\label{sec:revenue}

Let us first understand how the sellers' revenues depend on the prices $(p, q)$.

For $j \in \{1, 2\}$, let $r^-_j(p, q)$ be seller $j$'s revenue if the poor buyers
are forced to first buy as much of the other seller's dataset as possible,
and $r^+_j(p, q)$ be her revenue when the poor buyers are forced to buy
as much of seller $j$'s dataset as possible before buying from the other seller.

For notational convenience, denote the number of poor buyers by $n_p \defeq n-1$.

\begin{observation}
\label{thm:r-lo-hi}
For any $p, q \in \mathbb{R}_{\ge 0}$, we have
\begin{align*}
r^-_1(p, q) &= \min(p, \beta) + n_p \cdot \min(p, \max(0, 1-q))
& r^+_1(p, q) &= \min(p, \beta) + n_p \cdot \min(p, 1)
\\ r^-_2(p, q) &= n_p \cdot \min(q, \max(0, 1-p))
& r^+_2(p, q) &= n_p \cdot \min(q, 1)
\end{align*}
\end{observation}

\begin{observation}
For any $p, q \in \mathbb{R}_{\ge 0}$, we have
\begin{enumerate}
\item If $p + q \le 1$, then $r^-_j(p, q) = r^+_j(p, q)$ for all $j \in \{1, 2\}$.
\item If $p + q \ge 1$, then $r^-_1(p, q) + r^+_2(p, q) = r^+_1(p, q) + r^-_2(p, q) = \min(p, \beta) + n_p$.
\end{enumerate}
\end{observation}

Let $V(p, q) \subseteq \mathbb{R}_{\ge 0}^2$ be the set of \emph{valid} revenues of the sellers
when the datasets are priced at $p$ and $q$, respectively.
These are revenues obtainable when buyers behave as described in \cref{sec:intro}
(instead of being forced to purchase in a certain way.)

\begin{observation}
\label{thm:r}
For any $p, q \in \mathbb{R}_{\ge 0}$, if $(r_1, r_2) \in V(p, q)$, then
\begin{tightenum}
\item If $p + q \le 1$, then $r_j = r^-_j(p, q)$ for all $j \in \{1, 2\}$.
\item If $p < \alpha q$, then $r_1 = r^+_1(p, q)$ and $r_2 = r^-_2(p, q)$.
\item If $p > \alpha q$, then $r_1 = r^-_1(p, q)$ and $r_2 = r^+_2(p, q)$.
\item If $p = \alpha q$ and $p + q > 1$, then $r^-_1(p, q) \le r_1 \le r^+_1(p, q)$,
    $r^-(p, q) \le r_2 \le r^+_2(p, q)$, and $r_1 + r_2 = \min(p, \beta) + n_p$.
\end{tightenum}
\end{observation}

We now look at the maximum revenue a seller can obtain by setting a price strategically
given the other seller's price. Formally,
\begin{align*}
r_1^*(q) &\defeq \sup\{r_1: p \in \mathbb{R}_{\ge 0} \text{ and } (r_1, r_2) \in V(p, q)\}
\\ r_2^*(p) &\defeq \sup\{r_2: q \in \mathbb{R}_{\ge 0} \text{ and } (r_1, r_2) \in V(p, q)\}
\end{align*}

\begin{lemma}
\label{thm:rstar}
For any $p, q \in \mathbb{R}_{\ge 0}^2$,
\begin{align*}
r_1^*(q) &= \max\left(
    \beta + n_p\max(0, 1-q),
    \min(\beta, \alpha q) + n_p\min(1, \alpha q)
    \right),
\\ r_2^*(p) &= n_p\max(1-p, \min(1, p/\alpha)).
\end{align*}
\end{lemma}
\begin{proof}[Proof sketch]
Seller 1's revenue is maximized by either setting the price to $\beta$,
or undercutting seller 2 by an infinitesimally small amount.
Seller 2's revenue is maximized by either setting the price to $1$,
or undercutting seller 1 by an infinitesimally small amount.
\end{proof}

The pair $(p, q)$ is a $c$-NE if for any $(r_1, r_2) \in V(p, q)$, we have
$r^*_1(q) \le c \cdot r_1$ and $r^*_2(p) \le c \cdot r_2$.
Note that $r^*_1(q) \ge \beta \ge 1$ and $r^*_2(p) > 0$.
We follow the convention that $x/y = \infty$ if $x > 0$ and $y = 0$.
Thus, $(p, q)$ is a $c$-NE iff
\[ \mu(p, q) \defeq \inf_{(r_1, r_2) \in V(p, q)} \max\left(\frac{r_1^*(q)}{r_1}, \frac{r_2^*(p)}{r_2}\right) \le c. \]
Thus, the stablest prices $(p, q)$ are those that minimize $\mu(p, q)$.

Our goal is to show instability, i.e., finding a sufficiently-large constant $c$
such that $\mu(p, q) \ge c$ for all $p, q \in \mathbb{R}_{\ge 0}^2$.
This would prove that for all $\eps > 0$, no $(c-\eps)$-NE exists.
